\section{来源层级与估值准入}

本报告建立四层来源体系。第一层是华为、上市公司公告、监管问答、政府/交易所及竞品官方产品页面；第二层是可下载、可全文提取的原始卖方 PDF；第三层是结构化行情和财务摘要，必须经第二数据源或公告校验；第四层是媒体、搜索摘要、投资者提问前提和市场传闻。第一层可以确认产品状态、客户、订单和历史财务；第二层可以提供盈利预测与外部估值锚；第三层只能支持快照和筛选；第四层只能生成待验证问题，不能进入核心结论或估值。

\begin{exhibitbox}[声明—证据—估值准入矩阵]
\begin{tabularx}{\textwidth}{L{3.2cm}L{3.4cm}X L{2.4cm}}
\toprule
\textbf{声明} & \textbf{最低证据} & \textbf{当前结论} & \textbf{估值处理}\\
\midrule
1,024 卡实物展示 & 华为官方现场/发布稿 & 已确认实物公开展示，性能为厂商口径 & 支持产业方向，不支持公司收入\\
8,192 卡已交付 & 客户验收或华为交付公告 & 未找到；当前为 2026Q4 最大设计 & 0 权重\\
256TB 为 HBM & 内存定义与物理配置 & 未找到；仅确认全局统一地址空间 & 禁止使用\\
公司为 Atlas 供应商 & 公司/华为点名平台与产品 & 未找到上游 A 股闭环 & 0 专属收入\\
公司有 Atlas 订单 & 订单、合同或收入披露 & 未找到 & 0 专属利润\\
华为/昇腾合作关系 & 公司年报或官方伙伴资料 & 申菱、拓维、讯飞存在不同层级关系 & 只支持观察标签\\
2026E/2027E 盈利预测 & 新鲜原始卖方 PDF & 12 个模型代码具正权重外部预测 & 可进入更广义业务估值\\
\bottomrule
\end{tabularx}
\sourcenote{来源注册表、逐声明审计与估值复算审计；详见附录“来源与声明审计”。}
\end{exhibitbox}

\section{事实校准：标题里最容易误读的四点}

\textbf{“第一个超节点”不准确。}华为在 Atlas 950 之前已有 Atlas 900 A3 SuperPoD 和 CloudMatrix384。2026 年 7 月的准确表述是“Atlas 950 物理系统首次公开展示”，而不是“华为第一次做出超节点”。前代 384 卡系统的商业部署证明华为具备超节点产品化经验，但不能外推新一代供应商沿用、单价或出货量。

\textbf{“1,024 卡”与“8,192 卡”不能互换。}1,024 卡是本次公开展示的系统边界，对应约 1EFLOPS FP8 和 2EFLOPS FP4；8,192 卡是最大设计，包含 128 个计算柜和 32 个互联柜，对应约 8/16EFLOPS。前者验证硬件存在，后者仍需整机、稳定性、软件利用率和客户验收验证。

\textbf{“256TB”不是 HBM 结论。}1,024 卡系统披露的是 256TB 全局统一地址空间，平均约 256GB/卡；950DT 路线图则是 144GB HiZQ 2.0 HBM/卡。二者口径不同，可能包含 CPU 或其他池化内存，也可能对应不同配置。华为未解释前，报告不得将 256TB 标为 HBM，也不得据此反推 HBM 供应量。

\textbf{“全光”不等于某家光模块公司订单。}华为披露多机柜全光 UB-Mesh，但没有公开端口数、模块速率、可插拔/CPO/LPO/OCS 路线、模块数量和供应商。华为还表示 UBoE 相对 RoCE 可能减少交换机和光模块数量，因此简单以卡数乘固定光模块用量会同时忽略拓扑和技术替代风险。

\section{数据覆盖与校验结果}

财务包覆盖 22 家公司 FY2025 和 2026Q1 实际数据；9 家披露 2026H1 预告，其余 13 家明确标注为未披露。市场包覆盖 22 家 2026-07-17 收盘价、总市值、流通市值代理、五日均量和 2026-06-30 最新官方北向持股。22 家腾讯与 AStock/Sina 收盘价全部一致；科创板成交量的“股/手”差异已经修正，避免 100 倍单位错误。

卖方样本包含 38 份可下载原始 PDF，重点覆盖华工科技、英维克、申菱环境、科华数据、沃尔核材、科大讯飞、航天电器、AI PCB/CCL 和行业报告。拓维信息在 180 天新鲜窗口没有找到原始覆盖，旧报告只保留为生态历史证据，估值权重为 0。所有外部目标价缺失项均写为“未披露”，没有把当前价 PE 表反推成券商目标价。

\begin{sourcequalitybox}[数据质量边界]
华为性能数字为厂商自述，没有统一精度、稀疏性、功耗、软件版本和独立同负载基准；上市公司 FY2025/2026Q1 结构化摘要为中等质量，重点公司已用年报文本交叉核验；2026H1 为未经审计预告；行情为 2026-07-17 收盘快照；北向持股为季度末数据。以上时点和口径均不得混用。
\end{sourcequalitybox}

\section{负面证据同样重要}

航天电器问答确认高速背板、液冷连接器与管路能力，但没有确认投资者问题中预设的“昇腾 950 关系”；强一股份没有确认昇腾 950 探针卡或 HBM 客户；四川长虹明确否认所称 AI 服务器 ODM 业务。三类结果分别被归为“能力存在但关系未确认”“问题前提无法验证”“明确排除”。负面证据不会从附件中消失，而会直接降低分类与估值权重。
